\subsection{分次环与分次模}

定义2. 给定一个环 A 和加法群 A 上的一个分次 \((\mathbf{A}_{\lambda})\) （类型为 \(\Delta\) 的），若该分次与 A 上的环结构相容，则是指

\[
\mathrm{A} _ {\lambda} \mathrm{A} _ {\mu} \subset \mathrm{A} _ {\lambda + \mu} \quad f o r a l l \lambda , \mu i n \Delta .\tag{1}
\]

带有此分次的环 A 这时称为一个类型为 \(\Delta\) 的分次环。

命题1. 若 \(\Delta\) 的每个元素都可消去，且 \((\mathbf{A}_{\lambda})\) 是一个类型为 \(\Delta\) 的、与环 \(\mathbf{A}\) 的结构相容的分次，则 \(\mathbf{A}_0\) 是 \(\mathbf{A}\) 的一个子环（特别地 \(1 \in \mathbf{A}_0\) ）。

由于按定义 \(\mathbf{A}_0\mathbf{A}_0\subset \mathbf{A}_0\) ，只需证明 \(1\in \mathbf{A}_0\) 。设 \(1 = \sum_{\lambda \in \Delta}e_{\lambda}\) 是 1 的齐次分量分解。若 \(x\in \mathbf{A}_{\mu}\) ，则 \(x = x.1 = \sum_{\lambda \in \Delta}xe_{\lambda}\) ；比较次数为 \(\mu\) 的分量（由于 \(\mu +\lambda = \mu\) 蕴含 \(\lambda = 0\) ），得 \(x = xe_0\) 。由于这个关系对 A 的每个齐次元素都成立，故它对所有 \(x\in \mathbf{A}\) 成立；特别地， \(1 = 1.e_0 = e_0\in \mathbf{A}_0\) 。

定义3. 设 A 是一个类型为 \(\Delta\) 的分次环， \((\mathbf{A}_{\lambda})\) 是它的分次，M 是一个左（相应地，右）A-模；加法群 M 上的一个分次 \((\mathbf{M}_{\lambda})\) （类型为 \(\lambda\) 的）称为与 M 上的 A-模结构相容，如果

\[
\mathbf {A} _ {\lambda} \mathbf {M} _ {\mu} \subset \mathbf {M} _ {\lambda + \mu} \quad (\text { resp. } \mathbf {M} _ {\mu} \mathbf {A} _ {\lambda} \subset \mathbf {M} _ {\lambda + \mu})\tag{2}
\]

对所有 \(\lambda, \mu\) （在 \(\Delta\) 中）成立。带有此分次的模 \(\mathbf{M}\) 这时称为一个类型为 \(\Delta\) 的左（相应地，右）分次模（在分次环 \(\mathbf{A}\) 上）。

当 \(\Delta\) 的元素都可消去时，由 (2) 和命题 1 可知，诸 \(\mathbf{M}_{\lambda}\) 都是 \(\mathbf{A}_0\) -模。

显然，若 A 是一个类型为 \(\Delta\) 的分次环，则左 A-模 \(A_{s}\) （相应地，右 A-模 \(A_{d}\) ）是类型为 \(\Delta\) 的分次模。

例. (1) 在任意环 A 上，类型为 \(\Delta\) 的平凡分次与环结构相容。若 A 赋予平凡分次，则 A-模 M 上的一个分次 ( \(M_{\lambda}\) ) （类型为 \(\Delta\) 的）与 A-模结构相容的必要且充分条件是：诸 \(M_{\lambda}\) 都是 M 的子模。

\begin{enumerate}
\def\labelenumi{(\arabic{enumi})}
\setcounter{enumi}{1}
\tightlist
\item
  设 A 是一个类型为 \(\Delta\) 的分次环，M 是一个类型为 \(\Delta\) 的分次 A-模， \(\rho\) 是 \(\Delta\) 到一个以 0 为单位元的交换幺半群 \(\Delta'\) 的同态。则 A 是类型为 \(\Delta'\) 的分次环，M 是类型为 \(\Delta'\) 的分次模（对所考虑的类型为 \(\Delta'\) 的分次而言——它们由 \(\rho\) 从 A 和 M 上类型为 \(\Delta\) 的分次按第 1 号例 1 的程序导出）：这由关系式 \(\rho(\lambda + \mu) = \rho(\lambda) + \rho(\mu)\) 立即得出。
\end{enumerate}

特别地，若 \(\Delta = \Delta_{1} \times \Delta_{2}\) 是两个交换幺半群的积，则投影 \(pr_{1}\) 和 \(pr_{2}\) 是同态，而相应的分次正是由类型为 \(\Delta\) 的分次导出的部分分次（第 1 号，例 3）；因此这些部分分次与 A 的环结构以及 M 的模结构相容。

类似地，若 \(\Delta = \Delta_0^{(I)}\) （其中 \(\Delta_0\) 是一个以 0 为单位元的交换幺半群），则类型为 \(\Delta_0\) 的全分次（第 1 号，例 4）——它借助余对角同态由 A（相应地，M）上类型为 \(\Delta\) 的分次导出——与 A 上的环结构（相应地，与 M 上的模结构）相容。

\begin{enumerate}
\def\labelenumi{(\arabic{enumi})}
\setcounter{enumi}{2}
\tightlist
\item
  设 A 是一个类型为 \(\Delta\) 的分次环，M 是一个类型为 \(\Delta\) 的分次 A-模， \(\lambda_0\) 是 \(\Delta\) 的一个元素；对所有 \(\lambda \in \Delta\) ，令 \(M_{\lambda}' = M_{\lambda + \lambda_0}\) ，并令 \(M'\) 为 Z-模 \(\bigoplus_{\lambda \in \Delta} M_{\lambda}'\) 。由于 \(A_{\lambda}M_{\mu}' \subset M_{\lambda + \mu + \lambda_0} = M_{\lambda + \mu}'\) ，故 \(M'\) 是一个 A-模，且诸 \(M_{\lambda}'\) 在 M' 上构成一个与 M' 的 A-模结构相容的、类型为 \(\Delta\) 的分次；这样定义的类型为 \(\Delta\) 的分次 A-模 M' 称为把 M 的分次平移 \(\lambda_{0}\) 而得的分次 A-模，记作 \(\mathbf{M}(\lambda_{0})\) 。当 \(\Delta\) 是一个群时，分次 A-模 M' 的底层 A-模就等同于 M。
\end{enumerate}

*(4) 设 B 为交换环。一个未定元的多项式环 B{[}X{]} 是 N 型分次的，其分次由子群 BX \(^{n}\) (n ≥ 0) 给出（参见 III, § 2, 第 9 号及 IV）。*

*(5) 设 B 为交换环，E 为 B-模，Q 为 E 上的二次型，C(Q) 为 Q 的 Clifford 代数（参见 IX, §9）。子 B-模 C \(^{+}\) (Q) 与 C \(^{-}\) (Q) 在 C(Q) 上构成一个与 C(Q) 上的环结构相容的 Z/2Z 型分次。*

注. (1) 最常用的分次是 Z 型或 \(Z^{n}\) 型的分次；当我们谈到分次（相应地，双分次、三分次等等）模和环而不指明类型时，理解为指的是 Z 型（相应地， \(Z^{2}\) , \(Z^{3}\) 等等的类型）的分次；N 型的分次环（相应地，分次模）也称为具有正次数的分次环（相应地，分次模）。

\begin{enumerate}
\def\labelenumi{(\arabic{enumi})}
\setcounter{enumi}{1}
\tightlist
\item
  当 \(\mathbf{Z}\) 赋予平凡分次时，类型为 \(\Delta\) 的分次 \(\mathbf{Z}\) -模恰好就是定义 1（第 1 号）中的分次交换群（其次数集合是一个交换幺半群）。
\end{enumerate}

定义4. 设 \(\mathbf{A}, \mathbf{A}'\) 是两个同一类型 \(\Delta\) 的分次环， \((\mathbf{A}_{\lambda})\) , \((\mathbf{A}_{\lambda}')\) 是它们各自的分次。环同态 \(h: \mathbf{A} \to \mathbf{A}'\) 称为分次的，如果 \(h(\mathbf{A}_{\lambda}) \subset \mathbf{A}_{\lambda}'\) （对一切 \(\lambda \in \Delta\) ）。

设 \(\mathbf{M}, \mathbf{M}'\) 是一个类型为 \(\Delta\) 的分次环上的两个类型为 \(\Delta\) 的分次模。设 \(u: \mathbf{M} \to \mathbf{M}'\) 是一个 A-同态， \(\delta\) 是 \(\Delta\) 的一个元素； \(u\) 称为次数为 \(\delta\) 的分次同态，如果 \(u(\mathbf{M}_{\lambda}) \subset \mathbf{M}_{\lambda + \delta}\) （对一切 \(\lambda \in \Delta\) ）。

设 A 是一个类型为 \(\Delta\) 的分次环， \(A'\) 是一个类型为 \(\Delta'\) 的分次环，而 \(\rho: \Delta \to \Delta'\) 是一个同态。环同态 \(h: A \to A'\) 称为分次的，如果 \(h\) 是类型为 \(\Delta'\) 的分次环之间的分次同态，其中 A 赋予类型为 \(\Delta'\) 的分次，它由其类型为 \(\Delta\) 的分次借助 \(\rho\) （第 1 号，例 2）导出；这就是说， \(h(A_{\lambda}) \subset A_{\rho(\lambda)}'\) （对一切 \(\lambda \in \Delta\) ）。

一个 A-同态 \(u: \mathbf{M} \to \mathbf{M}'\) 称为分次的，如果存在 \(\delta \in \Delta\) 使得 \(u\) 是次数为 \(\delta\) 的分次同态。若 \(u \neq 0\) 且 \(\Delta\) 的每个元素都可消去，则次数 \(\delta\) （ \(u\) 的）这时被唯一确定。

若 \(h: \mathbf{A} \to \mathbf{A}'\) , \(h': \mathbf{A}' \to \mathbf{A}''\) 是类型为 \(\Delta\) 的分次环的两个分次同态，则 \(h' \circ h: \mathbf{A} \to \mathbf{A}''\) 也是；一个映射 \(h: \mathbf{A} \to \mathbf{A}'\) 是分次环同构的必要且充分条件是： \(h\) 是双射，且 \(h\) 与逆映射 \(h'\) 都是分次同态；为此也只需 \(h\) 是一个双射的分次同态。由此看出，分次同态可取作类型为 \(\Delta\) 的分次环结构物种的态射（集合论，IV, § 2, 第 1 号）。

类似地，若 \(u: \mathbf{M} \to \mathbf{M}'\) 和 \(u': \mathbf{M}' \to \mathbf{M}''\) 是类型为 \(\Delta\) 的分次 A-模的两个分次同态，其次数分别为 \(\delta\) 和 \(\delta', u' \circ u: \mathbf{M} \to \mathbf{M}''\) 是次数为 \(\delta + \delta'\) 的分次同态。若 \(\delta\) 有一个逆元 \(-\delta\) （在 \(\Delta\) 中）且 \(u: \mathbf{M} \to \mathbf{M}'\) 是一个次数为 \(\delta\) 的双射分次同态，则其逆映射 \(u': \mathbf{M}' \to \mathbf{M}\) 是一个次数为 \(-\delta\) 的双射分次同态。由此同上可知，次数为 0 的分次同态可取作类型为 \(\Delta\) 的分次 A-模物种的态射。但一个双射分次同态 \(u: \mathbf{M} \to \mathbf{N}\) ，若其次数 \(\neq 0\) ，在 \(\mathbf{M}\) 和 \(\mathbf{N}\) 非零且 \(\Delta\) 的元素可消去时并不是分次 A-模同构。

例. (6) 若 M 是一个分次 A-模，且 \(\mathbf{M}(\lambda_{0})\) 是由平移（第 2 号，例 3）得到的分次 A-模，则 \(\mathbf{M}(\lambda_{0})\) 到 M 中的、在每个 \(M_{\lambda+\lambda_{0}}\) 上都与典范单射一致的 Z-线性映射是次数为 \(\lambda_{0}\) 的分次同态（当 \(\Delta\) 是群时它是双射的）。

\begin{enumerate}
\def\labelenumi{(\arabic{enumi})}
\setcounter{enumi}{6}
\tightlist
\item
  若 \(a\) 是一个次数为 \(\delta\) 的、属于 A 的中心的齐次元素，则任一分次 A-模 M 的位似 \(x \mapsto ax\) 是次数为 \(\delta\) 的分次同态。
\end{enumerate}

注. (3) 分次 A-模 M 称为分次自由 A-模，如果存在由齐次元素组成的 M 的一个基 \((m_{\iota})_{\iota \in I}\) 。设情形如此且 \(\Delta\) 是交换群；令 \(\lambda_{i}\) 为 \(m_{i}\) 的次数，并对每个 \(\iota\) 考虑移位 A-模 \(\mathrm{A}(-\lambda_{i})\) （第 2 号，例 3）；若 \(e_{i}\) 表示 A 的元素 1，视为次数为 \(\lambda_{i}\) 的元素（在 \(\mathrm{A}(-\lambda_{i})\) 中），则 A-线性映射 \(u: \bigoplus_{\iota \in I} \mathrm{A}(-\lambda_{i}) \to \mathrm{M}\) （它满足 \(u(e_{\iota}) = m_{\iota}\) ，对所有 \(\iota\) ）是一个分次 A-模同构。

仍假设 \(\Delta\) 是交换群，现设 N 是一个分次 A-模， \((n_{\iota})_{\iota \in I}\) 是 N 的一个齐次生成元系，并设 \(n_{\iota}\) 的次数为 \(\mu_{\iota}\) 。则 A-线性映射 \(v: \bigoplus_{\iota \in I} A(-\mu_{\iota}) \to N\) （它满足 \(u(e_{\iota}) = n_{\iota}\) ，对所有 \(\iota\) ）是一个次数为 0 的满射分次 A-模同态。若 N 是有限生成的分次 A-模，则总存在 N 的有限齐次生成元系，从而存在上述类型的满射同态，其中 I 是有限的。

